% 核心结论：纠正过程图同列端口及注入方向，网络与模式卡片完全分开
% 草图：[真实几何与节点 | 条件、推导、结论]；17cm印宽，最小字号8pt。
% 图契约：双排错位端口：下排0列注入，上排N列抽出；图中的端口、几何、箭头和方程对应本题。
% 版面：沿用前三模块的几何图+推导分区；目标印宽17cm，简图至中等密度。
% 布局：左侧真实结构，右侧关键关系；真实电路边不作为架构流程箭头处理。
\documentclass[tikz,border=7pt]{standalone}
\usepackage{fontspec,xeCJK,amsmath,amssymb}
\IfFontExistsTF{Times New Roman}{\setmainfont{Times New Roman}}{\setmainfont{TeX Gyre Termes}}
\IfFontExistsTF{Heiti SC}{\setCJKmainfont[AutoFakeBold=1.5]{Heiti SC}}{\setCJKmainfont[AutoFakeBold=1.5]{FandolHei-Regular}}
\usetikzlibrary{arrows.meta,bending,calc,patterns,positioning,decorations.pathmorphing}
\definecolor{ink}{HTML}{243745}
\definecolor{blue}{HTML}{275DAB}
\definecolor{green}{HTML}{177C63}
\definecolor{red}{HTML}{BD4935}
\definecolor{orange}{HTML}{AA6C16}
\definecolor{purple}{HTML}{715393}
\tikzset{every node/.style={font=\fontsize{9.3}{12.5}\selectfont,text=ink},
arrow/.style={-{Stealth[length=5pt 1.5,width'=0pt 3]},shorten >=2pt,shorten <=1pt,line width=1pt},
arrow short/.style={-{Stealth[length=3pt 1.0,width'=0pt 2.5]},shorten >=1pt,shorten <=0pt,line width=.8pt},
axis/.style={arrow,draw=ink!60,line width=.6pt},
vec/.style={arrow,draw=blue},
wire/.style={draw=ink!75,line width=.8pt},
guide/.style={dashed,draw=ink!45,line width=.5pt},
card/.style={rounded corners=2pt,draw=blue!65,fill=blue!4,align=center,text width=6.5cm,inner xsep=9pt,inner ysep=7pt},
title/.style={font=\fontsize{10.5}{14}\selectfont\bfseries,text width=15cm,align=center},
label/.style={font=\fontsize{8.4}{11}\selectfont,align=center},
note/.style={label,text width=6.1cm}}
\begin{document}
\begin{tikzpicture}[x=1cm,y=1cm]
\path[use as bounding box] (-8,-3.5) rectangle(8,3.65);
\node[title] at (0,3.25) {M4-091\quad 双排错位端口：下排0列注入，上排N列抽出};
\begin{scope}[shift={(-4.05,0)}]
\draw[wire,draw=ink](-2.1,0.9)--(-1.6749999999999998,0.9) (-1.125,0.9)--(-0.7,0.9);\begin{scope}[shift={(-1.4,0.9)},rotate=0.0]\draw[draw=ink,fill=white,line width=.8pt](-0.275,-.085)rectangle(0.275,.085);\end{scope}
\draw[wire,draw=ink](-0.7,0.9)--(-0.275,0.9) (0.275,0.9)--(0.7,0.9);\begin{scope}[shift={(0.0,0.9)},rotate=0.0]\draw[draw=ink,fill=white,line width=.8pt](-0.275,-.085)rectangle(0.275,.085);\end{scope}
\draw[wire,draw=ink](0.7,0.9)--(1.125,0.9) (1.6749999999999998,0.9)--(2.1,0.9);\begin{scope}[shift={(1.4,0.9)},rotate=0.0]\draw[draw=ink,fill=white,line width=.8pt](-0.275,-.085)rectangle(0.275,.085);\end{scope}
\draw[guide](-2.1,0.9)--(-2.8,0.9)(2.1,0.9)--(2.8,0.9);\draw[wire,draw=ink](-2.1,-0.9)--(-1.6749999999999998,-0.9) (-1.125,-0.9)--(-0.7,-0.9);\begin{scope}[shift={(-1.4,-0.9)},rotate=0.0]\draw[draw=ink,fill=white,line width=.8pt](-0.275,-.085)rectangle(0.275,.085);\end{scope}
\draw[wire,draw=ink](-0.7,-0.9)--(-0.275,-0.9) (0.275,-0.9)--(0.7,-0.9);\begin{scope}[shift={(0.0,-0.9)},rotate=0.0]\draw[draw=ink,fill=white,line width=.8pt](-0.275,-.085)rectangle(0.275,.085);\end{scope}
\draw[wire,draw=ink](0.7,-0.9)--(1.125,-0.9) (1.6749999999999998,-0.9)--(2.1,-0.9);\begin{scope}[shift={(1.4,-0.9)},rotate=0.0]\draw[draw=ink,fill=white,line width=.8pt](-0.275,-.085)rectangle(0.275,.085);\end{scope}
\draw[guide](-2.1,-0.9)--(-2.8,-0.9)(2.1,-0.9)--(2.8,-0.9);\draw[wire,draw=ink](-2.1,0.9)--(-2.1,0.275) (-2.1,-0.275)--(-2.1,-0.9);\begin{scope}[shift={(-2.1,0.0)},rotate=-90.0]\draw[draw=ink,fill=white,line width=.8pt](-0.275,-.085)rectangle(0.275,.085);\end{scope}
\fill[ink] (-2.1,0.9)circle(2pt);
\fill[ink] (-2.1,-0.9)circle(2pt);
\draw[wire,draw=ink](-0.7,0.9)--(-0.7,0.275) (-0.7,-0.275)--(-0.7,-0.9);\begin{scope}[shift={(-0.7,0.0)},rotate=-90.0]\draw[draw=ink,fill=white,line width=.8pt](-0.275,-.085)rectangle(0.275,.085);\end{scope}
\fill[ink] (-0.7,0.9)circle(2pt);
\fill[ink] (-0.7,-0.9)circle(2pt);
\draw[wire,draw=ink](0.7,0.9)--(0.7,0.275) (0.7,-0.275)--(0.7,-0.9);\begin{scope}[shift={(0.7,0.0)},rotate=-90.0]\draw[draw=ink,fill=white,line width=.8pt](-0.275,-.085)rectangle(0.275,.085);\end{scope}
\fill[ink] (0.7,0.9)circle(2pt);
\fill[ink] (0.7,-0.9)circle(2pt);
\draw[wire,draw=ink](2.1,0.9)--(2.1,0.275) (2.1,-0.275)--(2.1,-0.9);\begin{scope}[shift={(2.1,0.0)},rotate=-90.0]\draw[draw=ink,fill=white,line width=.8pt](-0.275,-.085)rectangle(0.275,.085);\end{scope}
\fill[ink] (2.1,0.9)circle(2pt);
\fill[ink] (2.1,-0.9)circle(2pt);
\fill[green] (-2.1,-0.9)circle(2pt);\node[label]at(-1.78,-.58){$A$};
\fill[red] (0.7,0.9)circle(2pt);\node[below right=4pt] at(0.7,0.9){$B$};
\draw[vec,green](-2.1,-2.05)--(-2.1,-1.0)node[pos=0,left]{$I$};\draw[vec,red](0.7,1.0)--(0.7,2.25)node[pos=1,right]{$I$};\node[label]at(-2.1,2.55){$j=0$};\node[label]at(0.7,2.55){$j=N$};\node[note]at(0,-2.8){每条实线支路含一只r；左右两端都向无穷延伸};\end{scope}
\node[card,text width=6.5cm,draw=blue!65,fill=blue!4]at(3.9,1.8){$\rho=2-\sqrt3$};
\node[card,text width=6.5cm,draw=green!65,fill=green!4]at(3.9,0){同相输运：$R_0=Nr/2$};
\node[card,text width=6.5cm,draw=orange!65,fill=orange!4]at(3.9,-1.8){反相局域：$R_1=r(1+\rho^N)/(2\sqrt3)$};

\end{tikzpicture}\end{document}